\chapter{Nilpotent Holonomy and Completeness}\label{ch:nilpotent-hol}
\section{Nilpotent Holonomy}
\subsection{Nilpotent Groups}

In this section, we will see that the condition that the holonomy group of an
affine manifold $M$ be nilpotent will be useful in determining whether the
affine structure is complete. The main goal of this chapter is to prove the
Markus conjecture in this case. The conjecture is posed in
\cite{markus1963cosmological} and states that for a compact affine manifold,
completeness of the affine structure is equivalent to parallel volume. This
conjecture is yet to be solved in full generality. This chapter is follows
closely the paper \cite{nilpotent-holonomy} and many of the results we state
were first proved there.

\begin{definition}[Central Series]
    Let G be a group. Then a central series for $G$ is a finite collection of normal subgroups
    such that
    \begin{center}
        $1=G_0 \leq G_1 \leq \dots \leq G_s = G$
    \end{center}
    and
    \begin{center}
        $\dfrac{G_{i+1}}{G_i} \leq Z(\dfrac{G}{G_i})$
    \end{center}
    where $Z(\dfrac{G}{G_i})$ denotes the center.
\end{definition}

\begin{definition}[Nilpotent Group]
    A group $G$ is said to be nilpotent if it admits a central series.
\end{definition}
\begin{remark}
    The \textit{nilpotent class} of $G$ is the length of the shortest central series.
\end{remark}

\begin{eg}
    An important example of nilpotent groups are the upper triangular matrices. For example, the group
    $$
        H(R) = \left\{\begin{pmatrix}
            1 & a & b \\
            0 & 1 & c \\
            0 & 0 & 1
        \end{pmatrix}\right\}
    $$
    is called the Heisenberg group and is nilpotent of class 2. This group notably shows up in quantum mechanics.
\end{eg}
For nilpotent groups, we have the following results regarding their group
cohomology.

\begin{lemma}
    Let G be a nilpotent group and $E$ a $G$-module. If $H^0(G, E) = 0$ then
    also $H^i(G, E) = 0$ for all $i \geq 0$.
    \label{lemma:nilpotent-group-coho}
\end{lemma}
\begin{proof}
    The proof of this can be found in \cite{Hirsch1978}.
\end{proof}

A consequence of this is the following lemma.
\begin{lemma}
    If $G$ is a nilpotent group and $E$ is a $G$-module, then the following statements are equivalent.
    \label{lemma:nilpotent-group-equiv-hom-coho}
    \begin{enumerate}
        \item $H^0(G, E) = 0$
        \item $H^0(G, E^*) = 0$
        \item $H_0(G, E) = 0$
        \item $H_0(G, E^*) = 0$
    \end{enumerate}
\end{lemma}

\begin{proof}
    See \cite{nilpotent-holonomy}.
\end{proof}
\subsection{Unipotence and the Fitting submodule}
In the case that the affine holonomy group is nilpotent, we will be able to
show that completeness is equivalent to a property called unipotence of the
linear holonomy. We will now outline some consequences of unipotence.

\begin{definition}[Unipotent]
    A $G$-module $E$ is called unipotent if $(g-I)^n = 0$ for every $g \in G$
    where $n = \text{dim}(E)$.
\end{definition}
\begin{definition}
    An affine representation is unipotent if its linear part defines a unipotent
    $G$-module.
\end{definition}

Although $E$ may not be unipotent itself, we can consider the smallest
submodule (possibly the zero module) of $E$ which is unipotent. We shall denote
this submodule by $E_U$.

\begin{remark}
    The submodule $E_U$ is called the \textit{Fitting submodule} of $E$.
\end{remark}

Often, it will be useful to quotient out by the unipotent part of a $G$-module.
The following lemma gives some insight on why this is a useful idea.
\begin{lemma}\label{lemma:fitting-submodule}
    Let $E_U$ be the Fitting submodule of $E$. Then
    \begin{center}
        $H^0(G, E/E_U) = 0$
    \end{center}
\end{lemma}

\begin{proof}
    Take $E_U$ to be the Fitting submodule and assume $E_U \neq E$ otherwise the result is trivial.
    Suppose that $\dim E_U = r$ so that $(g-1)^r \cdot y = 0$ for all $y \in E_U$.

    As $H^0(G, E/E_U)$ is the set of stationary points of the $G$-module $E/E_U$ we
    can take $x + E_U \in H^0(G, E/E_U)$ not the identity so that
    \begin{center}
        $g\cdot (x + E_U) = x+ E_U$\\
        $ \implies (g - 1)\cdot x \in E_U$
    \end{center}
    for all $g \in G$.

    Now we claim that the $G$-submodule spanned by $x$ and $E_U$ is unipotent of
    dimension $r+1$. As $x \notin E_U$ this submodule is strictly bigger than
    $E_U$. Now, letting $y \in E_U$ we have that
    \begin{align*}
        (g-1)^{r+1}\cdot (x + y) & = (g-1)^{r} \cdot ((g-1)x + (g-1)y) \\
                                 & = 0
    \end{align*}
    since $(g-1)x + (g-1)y \in E_U$ and $\dim E_U = r$. Hence this a unipotent submodule strictly bigger than $E_U$ which contradicts maximality.
\end{proof}

The following splitting lemma will prove important in our characterization of
completeness.

\begin{lemma}
    If $G$ is a nilpotent group and $E$ a $G$-module with Fitting submodule $E_U$, then there exists a
    unique submodule $F$ such that $E = E_U \oplus F$.\label{lemma:fitting-unique-decomp}
\end{lemma}
\begin{proof}

    Let $\beta$ be the induced representation $\beta: G \to GL(E/E_U)$ on
    $GL(E/E_U)$. By Lemma \ref{lemma:fitting-submodule}, $H^0(G, E/E_U) = 0$.

    We will make use of the following claim which we resolve in Proposition
    \ref{prop:proofy}.
    \begin{claim}
        A submodule $F \subset E$ is complementary to $E_U$ e.g. $E = E_U \oplus F$ if and only if
        $F=T(E/E_U)$ where $T:E/E_U \to E$ is an equivariant linear map with $P \circ T = 1_{E/E_U}$ with $P$ the
        canonical projection map.
    \end{claim}

    Now let $S: E/E_U \to E$ be a linear map satisfying $P \circ S = 1_{E/E_U}$.
    Then any such $T$ as above can be uniquely written as $T = R + S$ with $R:
        E/E_U \to E_U$. It is $T$-equivariant if and only if

    \begin{center}
        $R + S = g \circ (R + S)\circ \beta(g)^{-1}$\\
        \begin{equation}\label{eq:fixed-point}
            \implies R = g\circ R \circ \beta(g)^{-1} + g\circ S \circ \beta(g)^{-1} - S.
        \end{equation}
    \end{center}
    If we can show that there is a unique such $T$, then we have shown that $E$ splits uniquely into $E = E_U \oplus F$.
    This will follow if we can show there is a unique $R$ satisfying the above.

    To do this we define a representation on $\Hom(E/E_U, E)$ by
    \begin{center}
        $\gamma : G \to \Hom(E/E_U, E)$\\
        $\gamma(g)(R) = g\circ R \circ \beta(g)^{-1}$
    \end{center}
    and we define a crossed homomorphism for $\gamma$ by $\mu(g) = g\circ S \circ \beta(g)^{-1} - S$. This
    defines a crossed homomorphism  as
    \begin{align*}
        \mu(gh) & = g \circ h\circ S \circ \beta(h)^{-1}\circ \beta(g)^{-1} -S \\
                & = \gamma(g)( h\circ S \circ \beta(h)^{-1}) - S               \\
                & = \gamma(g)(\mu(h) + S) - S                                  \\
                & = \gamma(g) (\mu(h)) + g\circ S\circ \beta(g)^{-1} -S        \\
                & = \gamma(g)(\mu(h)) + \mu(g).
    \end{align*}
    As $P$ is equivariant and $P \circ S = 1_{E/E_U}$ we get
    \begin{align*}
        P \circ \mu(g) & =  P \circ g \circ S \circ \beta(g)^{-1} - P \circ S \\
                       & =  \beta(g) \circ P  \circ S \circ \beta(g)^{-1} - 1 \\
                       & = \beta(g) \circ \beta(g)^{-1} - 1                   \\
                       & = 0.
    \end{align*}
    Now we can we write (\ref{eq:fixed-point}) as
    \begin{center}
        $R = \gamma(g)\circ R + \mu(g)$
    \end{center}
    which by Proposition \ref{prop:ary-point} is saying that $R$ is a fixed point of the affine representation
    induced by $\gamma$ and $\mu$. We wish now to show that this affine representation has a unique fixed point.
    This is equivalent by Lemma \ref{lemma:nilpotent-group-coho} and Lemma \ref{lemma:nilpotent-group-equiv-hom-coho}
    to $H^{0}(G, \Hom(E/E_U)) = 0$.

    To this end we suppose that $R: E/E_U \to E$ is fixed under all $\gamma(g)$ for
    $g\in G$. Then, letting $d=\dim(E_U)$ for any $d$ elements $g_1, \dots, g_d \in
        G$ we have
    \begin{center}
        $(1-g_1)\circ \dots \circ (1-g_d)|_{E_U} = 0$
    \end{center}
    as they act unipotently on $E_U$. It now follows by equivariance that
    \begin{align*}
        (1-g_1)\circ \dots \circ (1-g_d) \circ R & = (1-g_1)\circ \dots \circ (R-g_d \circ R)                             \\
                                                 & = (1-g_1)\circ \dots \circ (R-R \circ \beta(g_d))                      \\
                                                 & = (1-g_1)\circ \dots \circ R \circ(1- \beta(g_d))                      \\
                                                 & \implies R \circ (1 - \beta(g_1))\circ \dots \circ (1- \beta(g_d)) = 0
    \end{align*}
    so that $R$ vanishes on vectors of the form
    \begin{center}
        $(1 - \beta(g_1))\circ \dots \circ (1- \beta(g_d))x$
    \end{center}
    for $x \in E/E_U$. Thus if all vectors are in this span, then the only such $R$ fixed under the action of $\gamma$ is the
    zero map. As $1 - \beta(g)$ is invertible on $E/E_U$ (it acts without unipotent part), this is equivalent to
    showing that $E/E_U$ is spanned by vectors of the form $(1 - \beta(g))x$ with $x \in E/E_U$. By Proposition \ref{prop:group-zero-homol-cohom}, this is equivalent to
    $H_0(G, E/E_U)= 0$ which is true by Lemma \ref{lemma:fitting-submodule} completing the proof.
\end{proof}

This proves our result once we have cleared up that the stated claim is indeed
true which we show in the following proposition.
\begin{prop}\label{prop:proofy}
    A submodule $F \subset E$ is complementary to $E_U$ e.g. $E = E_U \oplus F$ if and only if
    $F=T(E/E_U)$ where $T:E/E_U \to E$ is an equivariant linear map with $P \circ T = 1_{E/E_U}$ with $P$ the
    canonical projection map.
\end{prop}
\begin{proof}
    Suppose $F=T(E/E_U)$, with $T$ being $G$-equivariant so that $F$ is a
    $G$-module. If $T(x+E_U) = e \in E$, then $P(e) = e+E_U = x+E_U \implies x \in
        E_U$ and so by linearity $e = 0$ and so $F \cap E_U = \{0\}$. Clearly, $E_U
        \oplus F \subseteq E$. For any $x \in E$, $x-T(x+E_U) \in E_U$. We can thus
    write $x = (x-T(x+E_U)) + T(x+E_U) \in E_U + F$.

    Conversely, suppose that $F$ is a $G$-submodule and $F \cap E_U = \{0\}$ and $F
        \oplus E_U = E$. Then its easy to see that $P|_F : F \to E/E_U$ is bijective.
    It is $G$-equivariant as $P$ is and $F$ is a $G$-submodule. We can thus define
    $T$ as $T = P^{-1}|_F$.
\end{proof}

An important consequence of this theorem is that for \emph{indecomposable
    affine representations} of a nilpotent group the affine representation is
always unipotent. We recall the definition of an indecomposable affine
representation now and then prove this result.

\begin{definition}[Decomposable affine representation]
    Let $E$ be a $G$-module. An affine representation $\alpha: G \to \aff(E)$ is said to be \emph{decomposable} if there exists a decomposition
    $E = E_1 \oplus F$ with $E_1 \neq E$ such that the decomposition is invariant under the linear part of $\alpha$ and some coset $x+ E_1$ is invariant under the entire affine holonomy $\alpha$ where $x \in E$.
\end{definition}
\begin{definition}[Indecomposable affine representation]
    An affine representation $\alpha: G \to \aff(E)$ is called \emph{indecomposable} if it is not decomposable.
\end{definition}
\begin{lemma}\label{lemma:indecomposable-coset-invariant}
    Let $E_U \subset E$ be the Fitting submodule of an affine representation $\alpha: G \to  \aff(E)$. If $G$ is nilpotent then
    there is a coset of $E_U$ that is invariant under $\alpha$.
\end{lemma}
\begin{proof}
    Let $\mu: G \to E$ be the translational component of the affine representation $\alpha = (\lambda,\mu)$. We check that
    $$
        \mu': G \xrightarrow{\mu} E \xrightarrow{\pi} E/E_U
    $$
    is a crossed homomorphism for $\lambda' = \pi \circ\lambda$. Letting $g,h \in G$ we compute
    \begin{align*}
        \mu'(gh) & = \pi \circ \mu(gh)                              \\
                 & = \pi \circ (\mu(g) + \lambda(g)\mu(h))          \\
                 & = \pi \circ \mu (g) + \pi \circ \lambda(g)\mu(h) \\
                 & = \mu'(g) + \lambda'(g)\mu'(h).
    \end{align*}
    Now, it follows that $\alpha' = (\lambda', \mu')$ is an affine representation of $G$ on $E/E_U$. Furthermore, the following diagram commutes
    \[
        \begin{tikzcd}
            E \arrow[r, "\pi"] \arrow[d, "\alpha(g)"']
            & E/E_U \arrow[d, "\alpha^{\prime}(g)"] \\
            E \arrow[r, "\pi"']
            & E/E_U
        \end{tikzcd}
    \]
    for all $g \in G$. Hence $\pi$ is a $G$-equivariant map with respect to the
    representations $\alpha$ and $\alpha'$. Now Lemma \ref{lemma:fitting-submodule}
    says that $H^0(G, E/E_U)=0$ and hence by Lemma \ref{lemma:nilpotent-group-coho}
    we have that $H^1(G, E/E_U) = 0$ since $G$ is nilpotent. Namely, the radiance
    obstruction $[c_{\alpha'}] \in H^1(G, E/E_U)$ vanishes and hence the affine
    representation $\alpha'$ fixes some $x + E_U$ where $x \in E$ and so $\alpha$
    fixes that coset.
\end{proof}

We can now prove the following theorem relating indecomposable affine
representations to unipotence.
\begin{lemma}\label{lemma:indecomposable-unipotent}
    Let $\alpha: G \to \aff(E)$ be an indecomposable affine representation of a nilpotent group $G$. Then $\alpha$ is unipotent.
\end{lemma}
\begin{proof}
    This follows as a corollary of our previous two lemmas. Lemma \ref{lemma:fitting-unique-decomp} tells us that there is a unique $F \subset E$
    such that $E = E_U \oplus F$ where $F$ is a $G$-submodule of $E$ and $E_U$ is the Fitting submodule. Now, Lemma \ref{lemma:indecomposable-coset-invariant} tells us
    that there is some coset of $E_U$ invariant under the linear part of $\alpha$. But this by definition means that $\alpha$ is decomposable unless $E_U = E$ which completes the proof as $E_U$
    is the unipotent part of $E$.
\end{proof}
\section{Unipotent Representations and Completeness}

\subsection{Expansions in the linear holonomy}
The aim of this section is to show that for a compact affine manifold, the
holonomy group being unipotent is equivalent to completeness of the affine
structure. Of geometrical importance to this will be the concept of an
\textit{expansion} in the linear holonomy.

\begin{definition}[Expansion]
    Let $T: V \to V$ be a linear map. Then $T$ is said to be an expansion of $V$ if for every eigenvalue
    $\lambda$ of $T$ we have $|\lambda|>1$.
\end{definition}

We now wish to prove Lemma \ref{lemma:exists-map} involving expansions in the
linear holonomy group. However, in order to do so we will need to recall some
of the theory of algebraic groups and the representation theory of nilpotent
groups.

\subsection{Algebraic groups and representation theory of Lie algebras}

We will need some tools from the theory of algebraic groups and the
representation theory of Lie algebras. This discussion will be extremely brief
and we refer the reader to \cite{humphreys1975linear} for more background on
algebraic groups and \cite{hall2003lie} for the representation theory of Lie
algebras.

We start with the definition of an algebraic group.
\begin{definition}[Algebraic Group]
    An \textit{algebraic group} is a matrix group that can be defined using polynomials.
\end{definition}

If $G$ is an algebraic group, then there is a unique irreducible component
passing through the identity denoted by $G^{\circ}$. The topology on $G$ is
taken to be the Zariski topology.

\begin{remark}
    This unique irreducible component is called the \textit{identity component} of $G$.
\end{remark}

\begin{prop}
    Let $G$ be an algebraic group. Then $G^{\circ}$ is a normal subgroup of finite index in $G$ whose cosets are connected and
    irreducible components of $G$. \label{prop:identity-comp}
\end{prop}

\begin{proof}
    See section $7.3$ of \cite{humphreys1975linear}.
\end{proof}

\begin{prop}
    Let $G$ be an algebraic group and $H$ a normal subgroup. Then
    \begin{center}
        $\overline{[G, H]} = [\overline{G}, \overline{H}]$
    \end{center}
    where $\overline{H}$ denotes the Zariski closure. \label{prop:zariski-closure}
\end{prop}
\begin{proof}
    This follows from Proposition 2.4 of \cite{borel1991linear}.
\end{proof}

This proposition leads to the following corollary.
\begin{corollary}\label{corr:closure-nilpotent}
    Let $G$ be an algebraic group and $N$ a normal subgroup that is nilpotent. Then $\overline{N}$ is also nilpotent.
\end{corollary}

\begin{proof}
    By Proposition \ref{prop:zariski-closure}, a central series for $N$ will pass to a central series of $\overline{N}$.
\end{proof}

The next result we will state is an important result from the representation
theory of Lie algebras. It will guarantee us a basis for our representation
such that the matrices of the representation are upper triangular.

\begin{definition}[Weight Space]
    Let $\mathfrak{g}$ be a Lie algebra and $\pi: \mathfrak{g} \to \mathfrak{gl}_V$ be a representation on a vector space $V$. For $\lambda \in \mathfrak{g}^*$
    we define the weight space of $\mathfrak{g}$ with respect to $\lambda$ to be
    \begin{center}
        $V_{\lambda}^{\mathfrak{g}} = \{v \in V | \pi(g)v = \lambda(g)v ,\quad \forall g \in \mathfrak{g}\}$.
    \end{center}
\end{definition}

\begin{remark}
    If $V_{\lambda}^{\mathfrak{g}} \neq 0$ we will say that $\lambda$ is a weight for $\pi$.
\end{remark}

\begin{lemma}[Lie's theorem]
    Let $\mathfrak{g}$ be a solvable Lie algebra and let $\pi$ be a representation of $\mathfrak{g}$ on a finite
    dimensional, non-zero vector space $V$. Then there exists
    a weight  $\lambda \in \mathfrak{g}^*$ for $\pi$.
\end{lemma}
\begin{proof}
    This result can be found for instance in \cite{kac2010lietheorem}.
\end{proof}

\begin{corollary}
    Let $V$ be a vector space. If $\pi : \mathfrak{g} \to \mathfrak{gl}(V)$
    is a finite dimensional representation of a solvable Lie algebra, then there is a basis of $V$ such that all linear transformations in
    $\pi(\mathfrak{g})$ are represented by upper triangular matrices.
\end{corollary}
We will need a slightly stronger theorem than Lie's theorem which gives a decomposition of our vector space
into generalized weight spaces.

\begin{definition}[Generalized Weight Space]
    Let $\pi$ be a representation of a Lie algebra $\mathfrak{g}$ over a finite dimensional vector space $V$. Then
    for $\lambda \in \mathfrak{g}^*$, we can associate to $\lambda$
    \begin{center}
        $V_{\lambda} = \{v \in V | (\pi(g) - \lambda(g))^nv =0 \text{ for all } g\in \mathfrak{g}\}$
    \end{center}
    where $n$ depends on $g \in \mathfrak{g}$ and $v \in V$. If $V_{\lambda} \neq 0$ then we call it
    the \textit{generalized weight space} associated to $\lambda$.
\end{definition}

\begin{remark}
    This tells us that $\pi(g) - \lambda(g)$ is a nilpotent operator so that $\pi(g) = \lambda(g) + N(g)$ with $N(g)$ nilpotent.
\end{remark}

\begin{lemma}
    \label{lemma:lie-algebra-decomp}
    Let $\mathfrak{g}$ be a nilpotent Lie algebra over $\mathbb{C}$. Let $\pi$ be a representation of
    $\mathfrak{g}$ on a finite dimensional complex vector space $V$. Then there are finitely many generalized
    weight spaces $ V_{\lambda_k}$, which are invariant under the action of the representation and $V = \oplus_k V_{\lambda_k}$.
\end{lemma}

\begin{proof}
    See Proposition 3.2 in \cite{ward2013semisimple}.
\end{proof}

We now have the required machinery to start our proof.

\begin{lemma}\label{lemma:exists-map}
    Let $G$ be a nilpotent group and $E=\R^n$ a $G$-module. Suppose that the linear holonomy group does not
    contain an expansion of $E$. Then for each $n \in \mathbb{N}_{>0}$, there is a $C^r$ map $\varphi : E \to \mathbb{R}$ for any $r \in \Z_{>0}$
    which satisfies the following.
    \begin{enumerate}
        \item $\varphi > 0$ almost everywhere.
        \item $\varphi$ is $G$-invariant.
        \item There exists some $a>0$ such that
              \begin{center}
                  $\varphi(e^{t}x) = e^{ta}\varphi(x)$
              \end{center}
              for all $t \in \mathbb{R}$ and $x \in E$.
    \end{enumerate}
\end{lemma}

\begin{proof}
    If we can prove this for a normal subgroup $G_0$ of finite index, then we will be able to construct our map $\varphi$ as follows.
    Let $g_1G_0, \dots g_kG_0$ be the left cosets of $G_0$ and $\varphi_0$ a map satisfying properties $1$ and $3$ defined on $G_0$.
    Now define
    $$\varphi(x) = \sum^{k}_{i=1} \varphi_0(g_ix).$$

    This will be $G$-invariant as $G_0$ is a normal subgroup of $G$ and
    multiplication by an element of $G$ will simply permute the cosets.

    As we want to eventually apply some results from complex Lie Theory, we
    complexify the dual space $E^*$ as $F = \mathbb{C}\otimes E^*$. The
    contragradient representation of $G$ on $E^*$ extends to a representation on
    $F$ in the obvious way. We will denote this as $\rho : G \to GL(F)$. Now,
    $\rho(G)$ is a nilpotent subgroup of $GL(F)$ that passes through the identity.

    Let $H$ be the identity component of the algebraic closure of $\rho(G)$. As
    $\rho(G)$ is nilpotent, by Corollary \ref{corr:closure-nilpotent} so is its
    Zariski closure and hence its identity component as subgroups of nilpotent
    groups are nilpotent. The subgroup $H$ will be a normal subgroup of finite
    index by Proposition \ref{prop:identity-comp} and so $G_0 = \rho^{-1}(H)$ will
    also be a normal subgroup of finite index.

    We can now obtain a Lie algebra representation of the Lie algebra of $H$ on $F$
    by
    \begin{center}
        $di: \mathfrak{h} \to \mathfrak{gl}(F)$
    \end{center}
    where $i$ is the inclusion map. Importantly, by construction on the underlying vector spaces $\rho(G_0) \subseteq di(\mathfrak{h})$. As $H$ is a nilpotent Lie group, its Lie algebra is nilpotent hence solvable so we can
    apply Lemma \ref{lemma:lie-algebra-decomp} to obtain a decomposition $F = \oplus^m_{k=1} F_k$ which is $di(\mathfrak{h})$-invariant and hence $\rho$-invariant by
    construction when restricted to $G_0$.

    Again, Lemma \ref{lemma:lie-algebra-decomp} yields that there is a basis $B_k$
    of each $F_k$ that represents the operators $di(h)|_{F_k}$ as upper triangular
    complex matrices
    \begin{center}
        $di_k(h) = \lambda_k(h)I + N_k(h)$
    \end{center}
    where $di_k$ is the restriction to $F_k$ and $N_k(h)$ is strictly upper triangular nilpotent matrix and $\lambda_k(h)$ are the eigenvalues of $\rho(h)$.
    The first basis vector in each $B_k$ is killed by the strictly upper triangular matrix so that if $f_k\in F_k$ is the
    first basis vector in $B_k$
    \begin{center}
        $di_k(h)f_k = \lambda_k(h)f_k$
    \end{center}
    for all $h \in H$.

    Now we define the following group morphisms
    \begin{center}
        $\varphi_k : G_0 \to \mathbb{R}$\\
        $ \varphi_k(g) = \log|\lambda_k(g)|$
    \end{center}
    for $k = 1,\dots, m$.

    This defines a map into $\mathbb{R}^m$ by
    \begin{center}
        $\varphi=(\varphi_1 , \dots , \varphi_m): G_0 \to \mathbb{R}^m.$
    \end{center}
    If $G_0$ does not contain any expansions of $E$ then the image $\varphi(G_0)$ is disjoint from the set
    $$
        P = \{y \in \mathbb{R}^m: y_k>0, k=1,\dots ,m\}.
    $$
    As it is a subgroup, it follows that its linear span $\text{span}\{\varphi(G_0)\}$ is also disjoint from $P$ as if a point in the image was
    in the strictly negative orthant, it's inverse would lie in $P$. Therefore there exists a vector in the orthogonal complement
    of $\varphi(G_0)$ that lies in the closure of $P$. We denote this by $v = (v_1,\dots, v_m)$ and it satisfies
    \begin{center}
        $v_k \geq 0$ for all $k$,
        \[\sum^m_{k=1}v_k = a \geq 0,\]
        \[\sum^m_{k=1}v_k\varphi_k(g) =0 \text{ for all } g\in G_0.\]
    \end{center}
    If we scale our $v$ by a positive scalar $c>0$ it will satisfy the same properties so that we can arrange that
    \begin{center}
        $v_k =0$ or $v_k>2r$ for each k with $r$ coming from the lemma.

        Now for each $k$ we have a vector $f_k\in \Hom(\mathbb{C}\otimes E,
            \mathbb{C})$ satisfy $di_k(h)f_k = \lambda_k(h)f_k$. We now include $E$ into
        $\mathbb{C}\otimes E$ and construct the map
        \begin{center}
            $\psi: E \to \mathbb{R}$
            \[\psi(x) = \prod_{k=1}^{m}|f_k(x)|^{v_k}.\]
        \end{center}
    \end{center}
    Now $\psi$ is $C^r$ differentiable as $|f_k(x)|^{v_k} = (u_k(x)^2 + w_k(x)^2)^m$ with $m>r$. Properties
    $1)$ and $3)$ are satisfied by our construction of $v_k$. We now show that it is $G_0$-invariant.
    \[\psi(g^{-1}x) = \prod_{k=1}^{m}|f_k(g^{-1}x)|^{v_k} \]
    noting this is the contragradient representation and also that $f_k$ is an
    eigenvector
    \begin{align*}
        \psi(g^{-1}x) & = \prod_{k=1}^{m}|g\circ f_k(x)|^{v_k}                             \\
                      & = \prod_{k=1}^{m}|\lambda(g)f_k(x)|^{v_k}                          \\
                      & = \prod_{k=1}^{m}|\lambda_k(g)|^{v_k}\prod_{k=1}^{m}|f_k(x)|^{v_k}
    \end{align*}

    Now, recalling that $ \varphi_k(g) = \log|\lambda_k(g)| \implies
        e^{\varphi_k(g)}= |\lambda_k(g)|$ so we obtained
    \begin{align*}
        \prod_{k=1}^{m}|\lambda_k(g)|^{v_k}\prod_{k=1}^{m}|f_k(x)|^{v_k} & = e^{\sum_{k=1}^{m}v_k\varphi_k(g)}\psi(x) \\
                                                                         & = \psi(x)
    \end{align*}
    completing the proof.
\end{proof}

\subsection{Consequences of expansions in linear holonomy}
The previous theorem leads to several important results regarding completeness
of affine manifolds.
\begin{theorem}
    Let $M$ be a compact affine manifold. Let $E_0 \subset E$ be a proper linear subspace invariant
    under the affine holonomy and denote by $\Lambda$ the linear holonomy. If the image of the linear holonomy group $\Lambda_1 \subseteq GL(E/E_0)$ is nilpotent,
    then some element of $\Lambda_1$ expands $E/E_0$. \label{thm:expansion-in-holonomy}
\end{theorem}

\begin{proof}

    Let $q: E \to E/E_0$ be the canonical projection. Denote by $R$ the radiant
    vector field on $E/E_0$. If $\{y_1,\dots , y_m\}$ are coordinates on $E/E_0$,
    then this appears as $R = \sum_i y_i \frac{\partial}{\partial y_{i}}$. Then in
    each affine chart there is a vector field $X_i$ that is $\pi$-related to $R$.
    As $M$ is compact, we can cover it in finitely many charts and using a
    partition of unity subordinate to that covering form the vector field
    \[X = \sum_{i}\mu_i X_i\]
    on $M$.

    Now suppose that $\Lambda_1$ does not contain any expansions of $E/E_0$. Then
    by Lemma \ref{lemma:exists-map}, there is a $\Lambda_1$-invariant $C^1$ map
    $\Psi: E/E_0 \to \mathbb{R}$. As the flow along the radiant vector field at a
    point $z$ is given by $\phi_t(z) = e^tz$, our map satisfies
    \begin{align*}
        d\Psi_z R(Z) & = \frac{d}{dt}\psi(e^tz)|_{t=0}     \\
                     & = \frac{d}{dt} e^{ta}\psi(z)|_{t=0} \\
                     & = a\Psi(z)
    \end{align*}
    for some $a>0$ and all $z\in E/E_0$.

    Now the composition
    \[\tilde{f}: \tilde{M} \to E \to E/E_0 \to \mathbb{R}\]
    is invariant under deck transformations. Indeed, let $\gamma \in \pi_1(M)$ then
    for $x \in \tilde{M}$

    \[\Psi\circ \pi \circ \text{dev}(\gamma \star x) = \Psi\circ \pi \circ \text{hol}(\gamma)\circ \text{dev}(x)\]

    and this is invariant as $E_0$ is invariant under the holonomy group.

    This implies that $\tilde{f}$ covers a map $f: M \to \mathbb{R}$. As by
    Proposition \ref{prop:dev-charts}, we can give charts in terms of the
    development pair, in affine coordinates we have that $f = \Psi \circ \pi$. Now

    \begin{align*}
        df_y X(y) & = d(\Psi \circ \pi)_y X(y)                      \\
                  & = \frac{d}{dt}\Psi \circ \pi (\gamma(t))|_{t=0} \\
                  & = \frac{d}{dt}\Psi(e^{t}\tilde{y})|_{t=0}       \\
                  & = a\Psi(\tilde{y})
    \end{align*}

    where $\tilde{y} = \pi(y)$ and $\gamma(t)$ is an integral curve of $X$ starting
    at $y$. This gives that

    \[df_y X(y) = af(y).\]

    It follows that if $\alpha: I \to M$ is an integral curve of $X$ starting at $p
        \in M$ then
    \begin{align*}
        f(\alpha(t)) & = \Psi\circ \pi (\alpha(t))  \\
                     & = \Psi(\overline{\alpha(t)}) \\
                     & = \Psi(e^t \tilde{p})        \\
                     & = e^{ta}\Psi(\tilde{p})      \\
                     & = e^{ta}f(\alpha(c)).
    \end{align*}

    Now as $\Psi > 0$ almost everywhere, there exists $x_0 \in M$ with $f(x_0)>0$
    and the integral curve $\alpha_0$ through $x_0$ is defined for all $t$ as M is
    compact. Then $\lim_{t \to \infty}f(\alpha_0 (t)) = \lim_{t \to \infty}
        e^{ta}f(\alpha_0(c)) \to \infty$. But as $f$ is a continuous function and $M$
    is compact it must be bounded and so we reach a contradiction.
\end{proof}

This proof leads to several important results. The first result follows almost
immediately.

\begin{corollary}\label{corr:radiant-expansion}
    Let $M$ be a compact radiant manifold with nilpotent linear holonomy $\Lambda$. Then $\Lambda$ contains an expansion of $E$.
\end{corollary}
\begin{proof}
    We can take $E_0 = 0$ in Theorem \ref{thm:expansion-in-holonomy} and the result follows immediately.
\end{proof}

The next corollary follows from considering the Fitting submodule of $E$.

\begin{corollary}
    Let $M$ be a compact affine manifold with nilpotent affine holonomy group. If $F \neq 0$ then some element of the linear holonomy group
    expands $F$.
    \label{corr:F-has-expansion}
\end{corollary}
\begin{proof}
    Recall from Lemma \ref{lemma:fitting-unique-decomp}, that nilpotence of the affine holonomy gives a unique decomposition
    $E = E_U \oplus F$ which is invariant under the action of the holonomy group. If we take $E_0 = E_U$ in Theorem \ref{thm:expansion-in-holonomy},
    then $F \cong E/E_U$ as a group module with respect to the linear holonomy group.
\end{proof}

\section{Parallel volume}

In \cite{markus1963cosmological} it is conjectured that parallel volume is
equivalent to completeness for affine manifolds. Using the previous results, we
can now give a partial resolution of this conjecture for the case that our
manifold has nilpotent holonomy. Orientable integral affine manifolds always
possess a parallel volume form since their linear holonomy is in
$\mathrm{SL}(\Z^2)$ and so the following results will prove important in our
classification of such structures.

\begin{theorem}(\cite{nilpotent-holonomy}) \label{thm:parallel-implies-unipotent}
    Let $M$ be a compact affine manifold with nilpotent affine holonomy group. Suppose that there is a parallel volume form on $M$.
    Then the linear holonomy is unipotent.
\end{theorem}

\begin{proof}
    Again we will exploit that there is a unique decomposition $E = E_U \oplus F$ by Lemma \ref{lemma:fitting-unique-decomp}. Now let $g$ be
    any element of the linear holonomy group. As there is a parallel volume form by Lemma \ref{lemma:pullback-form-linearmap}, $\det(g)=1$. But by the unique decomposition
    \[1 = \det(g) = \det(g|_{E_U}) \det(g|_F).\]
    But $g|_{E_U}$ is a unipotent operator and so has determinant 1. This forces
    $\det(g|_F)=1$ so that $F$ cannot contain an expansion. But now Corollary
    \ref{corr:F-has-expansion} forces $F=0$ and so $E = E_U$.
\end{proof}

\begin{prop}\label{prop:involutive-distr}
    Let $M$ be an affine manifold and suppose that $M$ is compact and has unipotent affine holonomy representation. Then there exists a collection of $n$ regular, involutive distributions
    $\mathfrak{F}_i$ on $M$, parallel one-forms $\omega_i$ on the leaves of $\mathfrak{F}_i$ vanishing on $\mathfrak{F}_{i-1}$ and vector fields $X_i \in \mathfrak{X}(M)$ such that
    \begin{enumerate}
        \item The $X_i$ are tangent to the distributions $\mathfrak{F}_i$.
        \item The one-forms $\omega_i(X_i) = 1$ on each leaf of $\mathfrak{F}_i$.
    \end{enumerate}
\end{prop}
\begin{proof}
    Let $\hol(\pi_1(M)) = \Gamma$ be the affine holonomy group and $\text{Lin}(\Gamma)$ the linear holonomy group.

    Now, suppose that the linear affine holonomy representation is unipotent. Then
    $\text{Lin}(\gamma)$ is upper-triangular with respect to some suitable basis
    and hence there exists a linear flag
    \begin{center}
        $0 = F_0 \subset F_1 \subset \dots \subset F_n = \mathbb{R}^n$
    \end{center}
    preserved by the linear affine holonomy group. Furthermore, the induced action on the quotients $F_k/F_{k-1}$ is trivial and so
    the restriction to $F_k$ preserves a linear functional $l_k : F_k \to \mathbb{R}$ with kernel $F_{k-1}$ as shown in Section 17.5 of \cite{humphreys2012linear}  .

    This determines a family of involutive distributions $\mathfrak{F}_k \subset
        TM$ as follows. The differential of the developing map $d \dev_p$ is an
    isomorphism for each $p \in \tilde{M}$. We define $\tilde{\mathfrak{F}_i}$ by
    $$ \tilde{\mathfrak{F}_i} = d\dev^{-1}(F_i). $$ Now, by equivariance of the
    developing map we have
    \begin{align*}
        d \dev \circ d\gamma (\tilde{\mathfrak{F}_i}) & = d\hol(\gamma) \circ d \dev(\tilde{\mathfrak{F}_i})            \\
                                                      & = \text{Lin}(\hol(\gamma)) \circ d \dev(\tilde{\mathfrak{F}_i}) \\
                                                      & =  d \dev(\tilde{\mathfrak{F}_i})
    \end{align*}
    as $F_i$ is invariant under the linear holonomy. Hence the $\tilde{\mathfrak{F}_i}$ descend to distributions $\mathfrak{F}_i$ on $M$ as they are invariant
    under deck transformations. Since $\dev$ is a local diffeomorphism there exists a cover by affine charts $\{U_i\}_{i \in I}$ such that the $\mathfrak{F}_i$ are constant distributions and hence they are integrable.

    Furthermore, each linear map $l_i$ determines a parallel one-form $\omega_i$
    defined on $\mathfrak{F}_i$ and vanishing on $\mathfrak{F}_{i-1}$. Indeed, we
    define these on $\tilde{M}$ by $$ \tilde{\omega_i}|_p(v) = l_i \circ d
        \dev_p(v) $$ for $v \in \tilde{\mathfrak{F}_i}|_p$. Since the $l_i$ are
    invariant under the linear holonomy, the $\tilde{\omega_i}$ descend to
    one-forms $\omega_i$ on $M$ which are parallel as they are constant tensors in
    affine charts where $\dev$ is a diffeomorphism.

    We can construct vector fields $X_i$ on M, tangent to $\mathfrak{F}_i$ such
    that $\omega_i(X_i)=1$ using partitions of unity. Locally there exist affine
    charts $\{U_k\}_{k \in I}$ such that the distributions $\mathfrak{F}_i$ are
    constant as they are the pull back of a constant distribution on $\R^n$ by an
    isomorphism. In each of these charts $U_k$, we can find a vector $v \in
        \mathfrak{F}_i|_p$ such that $\omega_i|_p(v) = 1$ where $p \in U_k$. We can
    extend this to a vector field $X_k$ in each chart such that $\omega_i(X_k(p)) =
        1$ on all of $U_k$.

    Now, we pick a partition of unity $\{\rho_k\}_{k \in I}$ subordinate to the
    cover $\{U_k\}_{k \in I}$ and construct a vector field $$ X_i = \sum_k \rho_k
        X_k. $$ This vector field is tangent to the $\mathfrak{F}_i$ and satisfies
    $\omega_i(X_i) = 1$ on the leaves of $\mathfrak{F}_i$. These lift to vector
    fields $\tilde{X}_i$ on the universal cover which are complete as $M$ is
    compact.
\end{proof}

This will be used in conjunction with the previous theorem in order to show
that for a compact affine manifold, unipotence of the affine holonomy
representation implies completeness of the affine structure.

\begin{theorem}\label{thm:unipotent-complete}
    Let $M$ be an affine manifold. If $M$ is compact and has unipotent linear holonomy representation then $M$ is complete.
\end{theorem}

\begin{proof}
    We pick a basepoint $x_0$ in $M$. Let $p: \tilde{M} \to M$ be the universal cover of $M$ with $p(\tilde{x_0}) = x_0$.
    We now choose a developing map $\dev: \tilde{M} \to \mathbb{R}^n$ such that $\text{dev}(\tilde{x_0}) = 0$. Let
    \begin{center}
        $0 = F_0 \subset F_1 \subset \dots \subset F_n = \mathbb{R}^n$
    \end{center}
    be the flag preserved by the linear holonomy and $\mathfrak{F}_i$ the corresponding involutive distribution. Let $\omega_i$ and
    $X_i$ be the corresponding one-forms and vector fields on $M$ satisfying the properties from Proposition \ref{prop:involutive-distr}. The $X_i$ lift to vector fields
    $\tilde{X}_i$ on $\tilde{M}$ which are complete as $M$ is compact and the $\omega_i$ lift to one forms $\tilde{\omega}_i$ on $\tilde{M}$ with $\tilde{\omega}(\tilde{X}_i) = 1$ on the leaves of
    $\tilde{\mathfrak{F}_i}$. We note that this is clear from the construction of these objects in Proposition \ref{prop:involutive-distr}.

    We first show surjectivity of the developing map. Let $\gamma_i(t)$ be an
    integral curve of $\tilde{X}_i$ starting at the point $\tilde{x_i}$ and let
    $\dev(\tilde{x_i}) = v_i \in \R^n$. From integrating $$
        \dfrac{d}{dt}\text{dev}\circ \gamma_i(t) = d\dev_{\gamma_i(t)} \circ
        \tilde{X}(\gamma_i(t)) $$ from $0$ to $t$ we obtain
    \begin{align*}
        \text{dev}\circ \gamma_i(t) & = \text{dev}\circ \gamma_i(0) + \int_{0}^{t}d\dev \circ \tilde{X}_i(\gamma_i(s))ds.
    \end{align*}
    Applying our linear functional $l_i$ to this we obtain
    \begin{align*}
        l_i(\text{dev}\circ \gamma_i(t)) & = l_i(\text{dev}\circ \gamma_i(0)) + \int_{0}^{t}l_i \circ d\dev \circ \tilde{X}_i(\gamma_i(s))ds \\
                                         & = l_i(\text{dev}\circ \gamma_i(0)) + \int_{0}^{t}\tilde{\omega_i}(\tilde{X}_i)_{\gamma(s)}ds      \\
                                         & = l_i(\text{dev}\circ \gamma_i(0)) + \int_{0}^{t}ds                                               \\
                                         & = l_i(\text{dev}\circ \gamma_i(0)) + t \tag{$\star$}. \label{eq:flow}                             \\
    \end{align*}

    Now we can start our inductive step. We start with some $v \in \R^n$ and we
    wish to show that there is a point $x \in \tilde{M}$ with $\dev(x) = v$. We
    will do so by flowing along the vector fields $\tilde{X}_i$ in a suitable
    manner.

    Taking $i=n$ we start at $\tilde{x}_0 \in \tilde{M}$ such that
    $\text{dev}(\tilde{x}_0) = 0$. We flow along $\tilde{X}_n$ for time $t=l_n(v)$
    until the point $\tilde{x}_1 \in \tilde{M}$ with $\text{dev}(\tilde{x}_1) =
        v_1$. From \eqref{eq:flow} we obtain

    \[l_n(v_1) = l_n(v) + 0\]
    and so $v- v_1 \in \ker l_n = F_{n-1}$.

    We can repeat the same process taking $i=n-1$ i.e. we flow along
    $\tilde{X}_{n-1}$ for $t=l_{n-1}(v-v_1)$ ending at $\tilde{x}_2 \in \tilde{M}$
    with $\text{dev}(\tilde{x}_2) = v_2$. Then
    \begin{align*}
        l_{n-1}(v_2) & = l_{n-1}(v-v_1) + l_{n-1}(v_1) \\
                     & = l_{n-1}(v)
    \end{align*}
    and hence we have that $v-v_2 \in \ker l_{n-1} = F_{n-2}$.

    We continue this process inductively until we reach $F_0 = \{0\}$ at which
    point we have $v-v_n = 0$ where $v_n = \dev(\tilde{x}_n)$. Hence we have showed
    surjectivity of the developing map.

    To show injectivity we consider a path
    \[\gamma_0 : [a,b] \to \tilde{M}\]

    which develops to a closed loop $\delta: [a,b] \to \mathbb{R}^n$. We can deform
    $\gamma_0$ to a new path $\gamma_1$ such that it develops entirely inside
    $F_{n-1}$. We define a path $\gamma_t : [a,b] \to \tilde{M}$ where
    $\gamma_t(s)$ is the image of $\gamma_0(s)$ at time $-tl_n(\delta(s))$ under
    the flow of $\tilde{X}_n$. Fix $s \in [a,b]$ and let $\alpha: [0,1]$ be an
    integral curve of $\tilde{X}_n$ with $\alpha(0) = \gamma_0(s)$. Then,
    \begin{align*}
        \alpha'(t) = \tilde{X}_n(\alpha(t))
    \end{align*}
    and applying the differential of the developing map to this yields
    \begin{align*}
        d\dev \circ \alpha' (t) = d \dev(\tilde{X}_n(\alpha(t))).
    \end{align*}
    We can apply our linear functional $l_n$ to this and obtain
    \begin{align*}
        l_n(d\dev \circ \alpha' (t) ) & = l_n(d \dev(\tilde{X}_n(\alpha(t))))      \\
                                      & = \tilde{\omega_i}(\tilde{X}_n(\alpha(t))) \\
                                      & = 1.
    \end{align*}
    Now, integrating this equation between $t=0$ and $t = -tl_n(\delta(s))$ yields
    \begin{align*}
        l_n(\dev \circ \gamma_t(s)) & = l_n(\dev \circ \gamma_0(s)) - tl_n(\delta(s)) \\
                                    & = l_n(\delta(s)) - tl_n(\delta(s))
    \end{align*}
    since $\dev \circ \gamma_0(s) = \delta(s)$ by definition of $\delta(s)$.

    At $t=1$ this yields that $\text{dev}\circ \gamma_t(s) \in \ker l_n = F_{n-1}$
    and this holds for all $s \in [a,b]$. Repeating this process inductively, we
    get a curve that develops entirely in $F_0 = \{0\}$. Since by construction, our
    original curve is path homotopic to this deformed curve that develops onto $0$,
    it follows as $\dev$ depends only on the homotopy class of the curve, that
    $\dev$ is injective.
\end{proof}

We summarize the previous results into the following theorem.

\begin{theorem}
    \label{thm:complete-parallel-volume}
    Let $M$ be a compact affine manifold with nilpotent affine holonomy representation. Then the following are equivalent
    \begin{enumerate}
        \item M is a complete affine manifold.
        \item The linear holonomy is unipotent.
        \item M has a parallel volume form.
        \item The affine holonomy is irreducible.
        \item The affine holonomy is indecomposable.
    \end{enumerate}
\end{theorem}
\begin{proof}
    We have that $2) \implies 1)$ by Theorem \ref{thm:unipotent-complete}. We have that $3) \iff 2)$ by Theorem \ref{thm:parallel-implies-unipotent}.
    Now, $1) \implies 4)$ by Theorem \ref{thm:complete-implies-irreducible} and by definition $4) \implies 5)$. Now $5) \implies 2)$ by Lemma \ref{lemma:indecomposable-unipotent} and we are done.
\end{proof}
\begin{remark}
    We note that if the fundamental group of $M$ is nilpotent, then its holonomy representation will also be
    nilpotent as this is preserved by group morphisms.
\end{remark}